\tituloPretextual{Resumo}

Este relatório apresenta o desenvolvimento de uma solução de automação para serviços de Administração de Sistemas em Linux, no âmbito da unidade curricular de Administração de Sistemas (2025/2026). O trabalho tem como objetivo reduzir tarefas manuais de configuração e promover maior consistência operacional através de scripts executados num servidor Linux virtualizado em VirtualBox. A solução contempla a automatização de operações sobre DNS, integração com serviços Web, gestão de partilhas de rede e procedimentos de apoio à administração, incluindo mecanismos de segurança e de backup. O projeto segue uma abordagem incremental, com validação funcional em ambiente de laboratório, recorrendo a uma máquina cliente para testes de conectividade e de acesso aos serviços disponibilizados. Para além da implementação técnica, o relatório sistematiza decisões de configuração, procedimentos de execução e critérios de verificação, constituindo um manual de referência para reprodução e manutenção do sistema.

\textbf{Palavras-chave}: administração de sistemas, Linux, automação, scripts, DNS, serviços de rede, segurança, backups.
